\documentclass[11pt]{article}
\usepackage[T1]{fontenc}
\usepackage{lmodern}
\usepackage[margin=1in]{geometry}
\usepackage{amsmath,amssymb,amsthm}
\usepackage{booktabs,microtype}
\usepackage[hidelinks]{hyperref}
\hypersetup{pdftitle={The associative-commutative spectrum of SC229},
  pdfauthor={Xiaojian Zeng}}
\newtheorem{theorem}{Theorem}
\newtheorem{lemma}[theorem]{Lemma}
\newcommand{\ind}{\mathbf{1}}
\DeclareMathOperator{\supp}{supp}
\title{The associative-commutative spectrum of SC229}
\author{Xiaojian Zeng\\[3pt]
  {\small PBC School of Finance, Tsinghua University, 100190}\\
  {\small\texttt{xiaojianzeng@foxmail.com}}}
\date{}

\begin{document}
\maketitle
\begin{abstract}
We determine the associative-commutative spectrum of the three-element
groupoid SC229. Its full linear term functions are in bijection with rooted
labelled trees having one distinguished child at every nonleaf vertex.
The key step reconstructs the tree from the unique Cartesian prime factors
of the function's truth set, with a separate treatment of singleton
indicators. Consequently, the exponential generating function satisfies
$T=x(1+Te^T)$, and the number of $n$-variable functions is
$c_n=(n-1)!\sum_{k=0}^{n-1}\binom{n}{k}k^{n-1-k}/(n-1-k)!$,
where $0^0=1$. We give canonical representatives, an integer recurrence,
and exact finite checks of the classification.
\end{abstract}

\section{The problem and its answer}

Let $A=\{0,1,2\}$, with the binary operation whose first argument indexes
the row in the following table:
\begin{equation}\label{eq:table}
\begin{array}{c|ccc}
\ast&0&1&2\\\hline
0&0&0&0\\
1&1&0&1\\
2&1&1&1
\end{array}
\end{equation}
For a finite nonempty set $I$ of labels, a \emph{full linear term on $I$}
is a binary term in which each variable $x_i$, $i\in I$, occurs exactly
once. Variable order and binary bracketing are arbitrary; constants are
excluded. Write $t^*:A^I\to A$ for its evaluation. With $[n]=\{1,\ldots,n\}$,
let $c_n$ count the distinct functions $t^*$ of full linear terms on $[n]$.
Labels are fixed coordinates, so we take no further quotient by
permutations. This is the associative-commutative spectrum in the sense of
Huang and Lehtonen~\cite{HuangLehtonen2023}.

Huang and Lehtonen~\cite[Section~7]{HuangLehtonen2024} identify
\eqref{eq:table} as SC229, give the initial values
$1,2,12,108,1340$, and ask for the general spectrum. They also record
\begin{equation}\label{eq:interchange}
((w\ast x)\ast y)\ast z=((w\ast x)\ast z)\ast y.
\end{equation}
The distinction between the first right factor and later right factors
suggested by this identity is exactly the information retained by a term
function.

A \emph{decorated tree on $I$} is a rooted tree with vertex set $I$ in
which each vertex having children has exactly one distinguished child.
There is no order on the other children. Every nonsingleton term has a
unique \emph{leftmost decomposition}
\begin{equation}\label{eq:leftmost}
t=(((x_i\ast R_1)\ast R_2)\cdots\ast R_m),\qquad m\geq1.
\end{equation}
It is obtained by following left operands down to a variable and recording
the right subterms from bottom to top. Define $D(x_i)$ to be the one-vertex
tree. Recursively, $D(t)$ has root $i$, distinguished child subtree $D(R_1)$,
and ordinary child subtrees $D(R_2),\ldots,D(R_m)$.

\begin{theorem}\label{thm:main}
For every finite nonempty label set $I$, two full linear terms $s,t$ on
$I$ satisfy $s^*=t^*$ if and only if $D(s)=D(t)$ as decorated labelled
trees. Every decorated tree on $I$ is realized by such a term. Thus
\begin{equation}\label{eq:sum}
c_n=(n-1)!\sum_{k=0}^{n-1}\binom{n}{k}
                 \frac{k^{n-1-k}}{(n-1-k)!}\qquad(n\geq1),
\end{equation}
with $0^0=1$. Equivalently, the formal exponential generating function
$T(x)=\sum_{n\geq1}c_nx^n/n!$ satisfies
\begin{equation}\label{eq:egf}
T=x(1+Te^T).
\end{equation}
\end{theorem}

The leftmost decomposition and its ordered rooted-tree encoding are due
to Huang and Lehtonen~\cite[Definitions~6.1.2--6.1.3]{HuangLehtonen2023}.
Their Proposition~6.1.4 permits forgetting all child order when
$(x\ast y)\ast z=(x\ast z)\ast y$. That hypothesis fails for SC229:
at $(x,y,z)=(2,0,1)$ the two sides are $0$ and $1$. The contribution here
is the injective reconstruction that retains precisely one distinguished
child, thereby proving that the tree model counts functions without
duplication. The tree encoding, Cartesian factorization, and coefficient
inversion are established tools.

Injectivity of a tree encoding is a reconstruction question. In a related
setting, Zeng~\cite[Theorem~1.1]{Zeng2026Multiplicity} reconstructs proper
trees from their chromatic symmetric functions when the weight of each
nonleaf core vertex occurs only once in the weighted core. Here each
coordinate label remains visible in the entire term function; the
reconstruction below uses this additional information to recover the
supports of the child subtrees. No chromatic-symmetric-function result is
needed for our proof.

\section{Evaluation and essential coordinates}

For every term $R$ put $q_R=\ind_{\{R^*=1\}}$. For a singleton,
$q_{x_j}(a)=\ind_{\{a_j=1\}}$, whereas $x_j^*(a)=a_j$ is a ternary
projection. Every nonsingleton term is $\{0,1\}$-valued by
\eqref{eq:table}, so $q_R=R^*$ in that case. We keep this distinction
throughout the recursive argument. Write $\supp(R)$ for its variable
label set.

For a Boolean function $q$ on coordinates disjoint from $i$, define
\[
H_i(q)(a)=\ind_{\{a_i=2\}}\vee
          \bigl(\ind_{\{a_i=1\}}\wedge(1-q(a))\bigr).
\]
The first multiplication in \eqref{eq:leftmost} gives $H_i(q_{R_1})$.
Every later left operand is a bit, and for $b\in\{0,1\}$ and $a\in A$
the table gives $b\ast a=b(1-\ind_{\{a=1\}})$. Consequently,
\begin{equation}\label{eq:evaluation}
q_t=H_i(q_{R_1})\prod_{j=2}^{m}(1-q_{R_j}).
\end{equation}
Products are ordinary products of bits, with empty product $1$.
The root coordinate and the child supports partition $I$.
Induction in \eqref{eq:evaluation} shows that $D(t)$ determines $q_t$.
Conversely, realize each child tree recursively, multiply its root
variable first by the distinguished child's term, and then by the other
child terms in any order. This realizes every decorated tree. It remains
to prove that the function determines the tree.

A coordinate is \emph{essential} if two assignments differing only there
give different function values.
\begin{lemma}\label{lem:essential}
For every full linear term $t$, the indicator $q_t$ is nonconstant,
vanishes at the all-zero assignment, and is essential in every coordinate
of $\supp(t)$. If $i$ is the leftmost label, then
$q_t|_{a_i=0}\equiv0$.
\end{lemma}
\begin{proof}
All assertions hold for $q_{x_i}=\ind_{\{a_i=1\}}$. Induct on support
size. For a nonsingleton, \eqref{eq:evaluation} gives the asserted zero
slice and the all-zero value. Set $a_i=2$ and all child coordinates to
zero; the value is $1$. Keeping those child coordinates fixed and changing
$a_i$ from $0$ to $2$ witnesses essentiality of $i$.

For a coordinate in $R_1$, use a pair of assignments witnessing its
essentiality in $q_{R_1}$, set $a_i=1$, and set all other child coordinates
to zero. The resulting function is $1-q_{R_1}$. For a coordinate in $R_j$,
$j\geq2$, use its essentiality witness in $q_{R_j}$, set $a_i=2$, and set
every other child to zero. The resulting function is $1-q_{R_j}$.
Each pair changes only the coordinate under consideration.
\end{proof}

\section{Disjoint prime factors}

A nonconstant Boolean function essential in all its coordinates is
\emph{AND-prime} if it cannot be written
$h(a)=b(a_U)c(a_V)$ for a partition of its support into two nonempty
sets $U,V$. Define \emph{OR-prime} using $b(a_U)\vee c(a_V)$ instead.
Primality here always concerns disjoint coordinate sets.

We use the standard uniqueness of Cartesian prime factorization of finite
relations; see Olteanu, Koch, and Antova~\cite[Section~6.1]{OlteanuKochAntova2008}.
The following argument includes the nonemptiness and essentiality details
needed here.

\begin{lemma}\label{lem:factor}
Let $h:A^I\to\{0,1\}$ be nonconstant and essential in every coordinate.
There is a unique partition $\mathcal P_*$ of $I$ and a unique AND-prime
Boolean factor $h_B$ on each $B\in\mathcal P_*$ such that
$h=\prod_{B\in\mathcal P_*}h_B$.
\end{lemma}
\begin{proof}
Let $S=h^{-1}(1)$, which is nonempty. Call a partition $\mathcal P$
factorizing if
$S=\prod_{B\in\mathcal P}\pi_B(S)$, where $\pi_B$ is coordinate
projection and products are taken on the indicated disjoint blocks.
Suppose $\mathcal P$ and $\mathcal Q$ both factorize $S$. Projecting
the $\mathcal Q$-product onto $B\in\mathcal P$ yields
\[
\pi_B(S)=\prod_{\substack{C\in\mathcal Q\\B\cap C\ne\varnothing}}
                       \pi_{B\cap C}(S).
\]
The choices on distinct blocks are independent; omitted blocks can be
filled because $S$ is nonempty. Substitution into the $\mathcal P$-product
shows that the common refinement also factorizes $S$.

There are finitely many partitions of $I$. The common refinement of all
factorizing partitions is therefore a factorizing partition $\mathcal P_*$,
finer than every other one. Each projection $\pi_B(S)$ is prime, since a
proper factorization of it would refine $\mathcal P_*$. Conversely, any
partition into prime projections must be $\mathcal P_*$: a proper
refinement within one of its blocks would contradict that block's
primality.

An AND factorization of $h$ gives just such a Cartesian factorization of
$S$. Each factor's truth set is its corresponding projection of $S$,
since the other nonempty factors can be filled. Thus the factors themselves
are uniquely determined. Essentiality of $h$ implies essentiality on each
factor's block, and excludes a constant-one factor on a nonempty block.
Nonemptiness excludes constant-zero factors. The prime projections hence
give exactly the asserted AND-prime factors.
\end{proof}

\begin{lemma}\label{lem:recognition}
For a nonsingleton term in \eqref{eq:leftmost}, the factors displayed in
\eqref{eq:evaluation} are precisely its disjoint AND-prime factors.
Among those factors, exactly one has an identically zero slice at value
$0$ of a coordinate. It is $H_i(q_{R_1})$, and that coordinate is $i$.
\end{lemma}
\begin{proof}
First, every $q_R$ is OR-prime. A singleton admits no proper coordinate
split. Otherwise, suppose $q_R=b(a_U)\vee c(a_V)$ on a proper split, and
put its leftmost label $r$ in $U$. Setting $a_r=0$ makes $q_R$ identically
zero by Lemma~\ref{lem:essential}, forcing $c\equiv0$. This makes all
coordinates in $V$ inessential, a contradiction. Boolean complementation
therefore makes $1-q_R$ AND-prime.

Put $J=\supp(R_1)$ and $H=H_i(q_{R_1})$. The function $H$ is essential
on $\{i\}\cup J$: its root-$0$ and root-$2$ slices distinguish $i$,
and its root-$1$ slice is $1-q_{R_1}$, essential on $J$. In an AND split
$H=b(a_U)c(a_V)$, put $i\in U$. The root-$2$ slice is identically one,
forcing $c\equiv1$. This contradicts essentiality on $V$. Thus $H$ is
AND-prime. The factors in \eqref{eq:evaluation} have disjoint nonempty
essential supports covering $I$, so Lemma~\ref{lem:factor} identifies
them as the unique prime factors.

The root-$0$ slice of $H$ is identically zero. No child coordinate has
that property in $H$, because assigning $a_i=2$ makes $H=1$ regardless
of that child's value. Every other factor $1-q_{R_j}$ is $1$ at its
all-zero assignment and therefore has no identically zero value-$0$
slice in any coordinate. This proves both uniqueness assertions. The
slice condition concerns each factor, rather than the entire product.
\end{proof}

\section{Reconstruction and normal forms}

We now reconstruct $D(t)$ from the function $q_t$ and its fixed label
set, inductively on the size of that set. A singleton indicator on $\{i\}$
is decoded as the leaf $i$. For $|I|\geq2$, Lemma~\ref{lem:factor}
recovers all prime factors from $q_t^{-1}(1)$, and
Lemma~\ref{lem:recognition} recovers the root factor $H$ and the root
coordinate $i$ using only those factors.

The distinguished child has support $J=\supp(H)\setminus\{i\}$ and
indicator
\begin{equation}\label{eq:recover}
q_{R_1}=1-H|_{a_i=1}.
\end{equation}
For every other prime factor $h_K$ on $K$, the corresponding ordinary
child indicator is $1-h_K$. Each recovered support is nonempty and
strictly smaller than $I$. Equation~\eqref{eq:evaluation} guarantees
that these recovered functions are term indicators. Induction recovers
their decorated trees, with \eqref{eq:recover} identifying the distinguished
one. No knowledge of the original syntax is used to choose a factor,
a coordinate, or a child support.

Hence equal indicators have equal decorated trees. For $|I|\geq2$ the
indicators are the term functions themselves, and for $|I|=1$ there is
only the single projection. Together with realization and the evaluation
formula, this proves the classification in Theorem~\ref{thm:main}.

It also gives explicit normal forms. On $[n]$, order the ordinary child
subtrees by their least vertex labels, which are distinct because the
supports are disjoint. Recursively choose their representatives and use
the distinguished child first in the leftmost product. Every function
has exactly one such representative. An arbitrary fixed total order
gives the same construction on any other finite label set; the underlying
tree classification does not depend on this choice.

\section{Counting the functions}

Let $B_r$ count sets of decorated trees whose supports partition a fixed
$r$-element label set, with $B_0=1$. Isolating the tree containing a
specified label gives the integer recurrence
\begin{equation}\label{eq:B}
B_r=\sum_{j=1}^{r}\binom{r-1}{j-1}c_jB_{r-j}\qquad(r\geq1).
\end{equation}
Choose the root of a nonleaf decorated tree, the $k$ labels of its
distinguished child subtree, that subtree, and the set of remaining
child subtrees. The classification gives
\begin{equation}\label{eq:c}
c_1=1,\qquad
c_n=n\sum_{k=1}^{n-1}\binom{n-1}{k}c_kB_{n-1-k}\qquad(n\geq2).
\end{equation}
These recurrences can be evaluated successively: after $c_1$, compute
$B_1,c_2,B_2,c_3,B_3,\ldots$. In particular, \eqref{eq:c} uses only
$c_k$ with $k<n$ and $B_r$ with $r\leq n-2$.

The exponential generating function for the sets counted by $B_r$ is
$e^T$: for each number of components, the labelled product counts ordered
components, and division by the factorial of their number forgets their
order. A decorated tree is a root with either no child or one distinguished
child and a set of other children. This gives \eqref{eq:egf} by the
usual labelled tree and forest construction; compare
Gessel~\cite[Section~3.2]{Gessel2016}.

For completeness, apply formal Lagrange inversion
\cite[Theorem~2.1.1]{Gessel2016} to $T=x\Phi(T)$, where
$\Phi(u)=1+ue^u$ and $\Phi(0)=1$. It yields
\begin{align*}
c_n
  &=(n-1)!\,[u^{n-1}](1+ue^u)^n\\
  &=(n-1)!\sum_{k=0}^{n-1}\binom{n}{k}
                     [u^{n-1-k}]e^{ku}\\
  &=(n-1)!\sum_{k=0}^{n-1}\binom{n}{k}
                     \frac{k^{n-1-k}}{(n-1-k)!}.
\end{align*}
For $n=1$, the convention $0^0=1$ gives $c_1=1$.
All manipulations are in formal power series over $\mathbb Q$; no
analytic convergence is needed. This completes Theorem~\ref{thm:main}.

\section{Exact finite checks and scope}

The accompanying standard-library Python program
\texttt{verify\_sc229.py} compares two enumerations. One generates every
binary bracketing and every coordinate permutation. The other generates
labelled trees by Pr\"ufer words, chooses every root, and chooses a
distinguished child at each nonleaf. The Pr\"ufer correspondence is
reviewed in Gessel~\cite[Section~4.6]{Gessel2016}. Both enumerations
evaluate representatives using all nine entries of \eqref{eq:table}.
Integer bitsets store the full truth tables, including all $3^n$
assignments per function.

\begin{center}
\begin{tabular}{r r r}
\toprule
$n$ & Full linear terms & Distinct functions and decorated trees\\
\midrule
1 & 1 & 1\\
2 & 2 & 2\\
3 & 12 & 12\\
4 & 120 & 108\\
5 & 1680 & 1340\\
6 & 30240 & 21150\\
\bottomrule
\end{tabular}
\end{center}

The two enumerations give identical function-to-tree mappings for
$1\leq n\leq6$, with no collisions. A separate truth-set decoder tests
binary coordinate splits for Cartesian factorization and reconstructs
all $1463$ distinct functions for $1\leq n\leq5$. It is not exhaustively
tested at $n=6$. The recurrence and finite sum agree in exact arithmetic
through $n=12$; values after $n=6$ in that check are formula evaluations,
not additional function censuses. The program also tests
\eqref{eq:interchange} on all $81$ four-variable assignments. It can be
run with Python~3 without third-party packages.

These computations corroborate the universal proof; the proof has no
computational premise. The result answers the full SC229 counting
question of Huang and Lehtonen. It does not classify other operations
satisfying \eqref{eq:interchange}, or terms with constants or repeated
variables. A bounded primary-source search completed on 12 September
2026 identified no equivalent SC229 classification in the sources
inspected; this is not an exhaustive priority assessment.

\paragraph{Computational and AI assistance.}
OpenAI Codex assisted with proof development, exact verification code,
source checking, and manuscript preparation. A separate model context
checked the mathematical argument and supplied the differently organized
finite enumeration underlying the accompanying program. Such a check is
not external peer review or statistically independent evidence. All
reported computations use exact integers, finite relations, and rational
arithmetic, without randomness. No proof assistant was used.

\bibliographystyle{plainurl}
\bibliography{references}
\end{document}
